\begin{frame}[fragile]{Laatste dingetjes} %shoutout ivo voor het stelen van de slides (basically)
Bij berekeningen wil je vaak vergelijkingen onder elkaar zetten. Je doet dit 
met de \mintinline{tex}{align} omgeving van \texttt{amsmath}-package. 
\begin{columns}[t]
    \begin{column}{0.5\textwidth}
    \begin{minted}{tex}
    Zo zet je het onder elkaar
    \begin{align}
        f(x) & = (x+1)(x-1) - (x+1)^2  \\
        & = (x^2 -1) - (x^2 + 2x + 1) \\
        & = -2x - 2
    \end{align}
    \end{minted}
    \par
\end{column}

\begin{column}{0.46\textwidth}
    \vspace{1mm}
    \cmfontblock{ \small

    \noindent
   Zo zet je het onder elkaar
    \begin{align}
        f(x) & = (x+1)(x-1) - (x+1)^2 \tag{1}\\
        & = (x^2 -1) - (x^2 + 2x + 1) \tag{2}\\
        & = -2x - 2 \tag{3}
    \end{align}
    }
\end{column}
\end{columns}
\end{frame}